% from main.tex
%%%%%%%%%%  scan idx 564  |  folio 551  %%%%%%%%%%
Donc il suffit de montrer que l'on a pour tout \(k>1\)
\[
 \det(\gamma^k(E-n))=1,
\]
où, comme ci-dessus, \(E=\cl^\circ(\underline E)\), \(n=\epsilon(E)\). Or, on a
\[
 \gamma^k(x)=\lambda^k(x+k-1)
\]
pour tout \(x\in K^\circ(X)\), de sorte que
\[
 \det(\gamma^k(E-n))
 =
 \det(\lambda^k(E-n+k-1))
 =
 \det\left(\sum_{p+q=k}\binom{n+k-1}{p}\lambda^q(E)\right).
\]
Mais (6.2.10) et (6.5.3) donnent
\[
 \det(\lambda^q(E))=(\det(E))^{\binom{n-1}{q-1}}
\]
pour tout \(q>0\), d'où
\[
 \det(\gamma^k(E-n))=(\det(E))^r,
\]
avec
\[
 r=\sum_{\substack{p+q=k\\ q\geq1}}
 \binom{-n+k-1}{p}\binom{n-1}{q-1}.
\]
Or, \(r\) est une section de \(\mathbb Z\), partout égale au coefficient de \(t^{k-1}\) dans \((1+t)^{k-2}\), donc nulle. On a gagné.

\section*{6. Appendice : Calcul des déterminants des faisceaux localement libres}

\subsection*{6.1}
Considérons d'abord un topos annelé quelconque \(X=(\mathcal C,\mathcal O)\). Pour toute section positive \(n\in\Gamma(\mathbb Z_X)\) du faisceau constant \(\mathbb Z\) et pour tout \(\mathcal O\)-module \(M\in\Mod_{\mathcal O}(X)\), il est clair comment définir les \(\mathcal O\)-modules
